\subsection{Contador Síncrono MOD 24}

\paragraph{Incremento}

A operação de incremento realiza a transição:
\[
D \mapsto (D + 1) \mod 24
\]


A Tabela abaixo descreve o comportamento completo da função de incremento:

\begin{center}
\begin{tabular}{cccccc|ccccccc}
$D11$&$D10$&$D03$&$D02$&$D01$&$D00$&$o11$&$o10$&$o03$&$o02$&$o01$&$o00$&$Carry$\\
\hline
$0$&$0$&$0$&$0$&$0$&$0$&$0$&$0$&$0$&$0$&$0$&$1$&$0$\\
$0$&$0$&$0$&$0$&$0$&$1$&$0$&$0$&$0$&$0$&$1$&$0$&$0$\\
$0$&$0$&$0$&$0$&$1$&$0$&$0$&$0$&$0$&$0$&$1$&$1$&$0$\\
$0$&$0$&$0$&$0$&$1$&$1$&$0$&$0$&$0$&$1$&$0$&$0$&$0$\\
$0$&$0$&$0$&$1$&$0$&$0$&$0$&$0$&$0$&$1$&$0$&$1$&$0$\\
$0$&$0$&$0$&$1$&$0$&$1$&$0$&$0$&$0$&$1$&$1$&$0$&$0$\\
$0$&$0$&$0$&$1$&$1$&$0$&$0$&$0$&$0$&$1$&$1$&$1$&$0$\\
$0$&$0$&$0$&$1$&$1$&$1$&$0$&$0$&$1$&$0$&$0$&$0$&$0$\\
$0$&$0$&$1$&$0$&$0$&$0$&$0$&$0$&$1$&$0$&$0$&$1$&$0$\\
$0$&$0$&$1$&$0$&$0$&$1$&$0$&$1$&$0$&$0$&$0$&$0$&$0$\\
$0$&$0$&$1$&$0$&$1$&$0$&$0$&$0$&$1$&$0$&$1$&$1$&$0$\\
$0$&$0$&$1$&$0$&$1$&$1$&$0$&$1$&$0$&$1$&$0$&$0$&$0$\\
$0$&$0$&$1$&$1$&$0$&$0$&$0$&$0$&$1$&$1$&$0$&$1$&$0$\\
$0$&$0$&$1$&$1$&$0$&$1$&$0$&$1$&$0$&$1$&$0$&$0$&$0$\\
$0$&$0$&$1$&$1$&$1$&$0$&$0$&$0$&$1$&$1$&$1$&$1$&$0$\\
$0$&$0$&$1$&$1$&$1$&$1$&$0$&$1$&$1$&$0$&$0$&$0$&$0$\\
$0$&$1$&$0$&$0$&$0$&$0$&$0$&$1$&$0$&$0$&$0$&$1$&$0$\\
$0$&$1$&$0$&$0$&$0$&$1$&$0$&$1$&$0$&$0$&$1$&$0$&$0$\\
$0$&$1$&$0$&$0$&$1$&$0$&$0$&$1$&$0$&$0$&$1$&$1$&$0$\\
$0$&$1$&$0$&$0$&$1$&$1$&$0$&$1$&$0$&$1$&$0$&$0$&$0$\\
$0$&$1$&$0$&$1$&$0$&$0$&$0$&$1$&$0$&$1$&$0$&$1$&$0$\\
$0$&$1$&$0$&$1$&$0$&$1$&$0$&$1$&$0$&$1$&$1$&$0$&$0$\\
$0$&$1$&$0$&$1$&$1$&$0$&$0$&$1$&$0$&$1$&$1$&$1$&$0$\\
$0$&$1$&$0$&$1$&$1$&$1$&$0$&$1$&$1$&$0$&$0$&$0$&$0$\\
$0$&$1$&$1$&$0$&$0$&$0$&$0$&$1$&$1$&$0$&$0$&$1$&$0$\\
$-$&$1$&$1$&$0$&$0$&$1$&$1$&$0$&$0$&$0$&$0$&$0$&$0$\\
$0$&$1$&$1$&$0$&$1$&$0$&$0$&$1$&$1$&$0$&$1$&$1$&$0$\\
$0$&$1$&$1$&$0$&$1$&$1$&$1$&$0$&$0$&$1$&$0$&$0$&$0$\\
$0$&$1$&$1$&$1$&$0$&$0$&$0$&$1$&$1$&$1$&$0$&$1$&$0$\\
$-$&$1$&$1$&$1$&$0$&$1$&$1$&$0$&$0$&$1$&$0$&$0$&$0$\\
$0$&$1$&$1$&$1$&$1$&$0$&$0$&$1$&$1$&$1$&$1$&$1$&$0$\\
$0$&$1$&$1$&$1$&$1$&$1$&$1$&$0$&$1$&$0$&$0$&$0$&$0$\\
$1$&$0$&$0$&$0$&$0$&$0$&$1$&$0$&$0$&$0$&$0$&$1$&$0$\\
$1$&$0$&$0$&$0$&$0$&$1$&$1$&$0$&$0$&$0$&$1$&$0$&$0$\\
$1$&$0$&$0$&$0$&$1$&$0$&$1$&$0$&$0$&$0$&$1$&$1$&$0$\\
$1$&$0$&$0$&$0$&$1$&$1$&$0$&$0$&$0$&$0$&$0$&$0$&$1$\\
$1$&$0$&$0$&$1$&$0$&$0$&$1$&$0$&$0$&$1$&$0$&$1$&$0$\\
$1$&$0$&$0$&$1$&$0$&$1$&$1$&$0$&$0$&$1$&$1$&$0$&$0$\\
$1$&$0$&$0$&$1$&$1$&$0$&$1$&$0$&$0$&$1$&$1$&$1$&$0$\\
$1$&$0$&$0$&$1$&$1$&$1$&$0$&$0$&$1$&$0$&$0$&$0$&$1$\\
$1$&$0$&$1$&$0$&$0$&$0$&$1$&$0$&$1$&$0$&$0$&$1$&$0$\\
$1$&$0$&$1$&$0$&$0$&$1$&$1$&$1$&$0$&$0$&$0$&$0$&$0$\\
$1$&$0$&$1$&$0$&$1$&$0$&$1$&$0$&$1$&$0$&$1$&$1$&$0$\\
$1$&$0$&$1$&$0$&$1$&$1$&$0$&$1$&$0$&$0$&$0$&$0$&$1$\\
$1$&$0$&$1$&$1$&$0$&$0$&$1$&$0$&$1$&$1$&$0$&$1$&$0$\\
$1$&$0$&$1$&$1$&$0$&$1$&$1$&$1$&$0$&$1$&$0$&$0$&$0$\\
$1$&$0$&$1$&$1$&$1$&$0$&$1$&$0$&$1$&$1$&$1$&$1$&$0$\\
$1$&$0$&$1$&$1$&$1$&$1$&$0$&$1$&$1$&$0$&$0$&$0$&$1$\\
$1$&$1$&$0$&$0$&$0$&$0$&$1$&$1$&$0$&$0$&$0$&$1$&$0$\\
$1$&$1$&$0$&$0$&$0$&$1$&$1$&$1$&$0$&$0$&$1$&$0$&$0$\\
$1$&$1$&$0$&$0$&$1$&$0$&$1$&$1$&$0$&$0$&$1$&$1$&$0$\\
$1$&$1$&$0$&$0$&$1$&$1$&$0$&$1$&$0$&$0$&$0$&$0$&$1$\\
$1$&$1$&$0$&$1$&$0$&$0$&$1$&$1$&$0$&$1$&$0$&$1$&$0$\\
$1$&$1$&$0$&$1$&$0$&$1$&$1$&$1$&$0$&$1$&$1$&$0$&$0$\\
$1$&$1$&$0$&$1$&$1$&$0$&$1$&$1$&$0$&$1$&$1$&$1$&$0$\\
$1$&$1$&$0$&$1$&$1$&$1$&$0$&$1$&$1$&$0$&$0$&$0$&$1$\\
$1$&$1$&$1$&$0$&$0$&$0$&$1$&$1$&$1$&$0$&$0$&$1$&$0$\\
$1$&$1$&$1$&$0$&$1$&$0$&$1$&$1$&$1$&$0$&$1$&$1$&$0$\\
$1$&$1$&$1$&$0$&$1$&$1$&$1$&$0$&$0$&$0$&$0$&$0$&$1$\\
$1$&$1$&$1$&$1$&$0$&$0$&$1$&$1$&$1$&$1$&$0$&$1$&$0$\\
$1$&$1$&$1$&$1$&$1$&$0$&$1$&$1$&$1$&$1$&$1$&$1$&$0$\\
$1$&$1$&$1$&$1$&$1$&$1$&$1$&$0$&$1$&$0$&$0$&$0$&$1$\\

\end{tabular}
\end{center}

A partir da tabela verdade, e utilizando mapas de Karnaugh, obtêm-se as seguintes expressões booleanas:

$o11 = D10 \cdot D03 \cdot D00+D11 \cdot  \overline{D01} +D11 \cdot  \overline{D00} $~\\
$o10 =  \overline{D10}  \cdot D03 \cdot D00+D10 \cdot  \overline{D03} +D10 \cdot  \overline{D00} $~\\
$o03 = D02 \cdot D01 \cdot D00+D03 \cdot  \overline{D00} $~\\
$o02 =  \overline{D11}  \cdot  \overline{D02}  \cdot D01 \cdot D00+D02 \cdot  \overline{D01} +D02 \cdot  \overline{D00} $~\\
$o01 =  \overline{D03}  \cdot  \overline{D01}  \cdot D00+D01 \cdot  \overline{D00} $~\\
$o00 =  \overline{D00} $~\\
$Carry = D11 \cdot D01 \cdot D00$~\\

O sinal $Carry$ indica a ocorrência de overflow, isto é, a transição de $101$ para $000$.

As Figuras~\ref{fig:incmod24.png} e~\ref{fig:incmod24-red.png} apresentam, respectivamente, a implementação lógica do circuito e sua realização em redstone.

\imgbig{incmod24.png}{Circuito lógico para incremento em módulo 24.}
\imgbig{incmod24-red.png}{Implementação em redstone do circuito de incremento em módulo 24.}

\paragraph{Decremento}

A operação de decremento realiza a transição:
\[
D \mapsto (D - 1) \mod 24
\]

A tabela correspondente é:

\begin{center}
\begin{tabular}{cccccc|ccccccc}
$D11$&$D10$&$D03$&$D02$&$D01$&$D00$&$Carry$&$o11$&$o10$&$o03$&$o02$&$o01$&$o00$\\
\hline
$0$&$0$&$0$&$0$&$0$&$0$&$1$&$1$&$0$&$0$&$0$&$1$&$1$\\
$0$&$0$&$0$&$0$&$0$&$1$&$0$&$0$&$0$&$0$&$0$&$0$&$0$\\
$0$&$0$&$0$&$0$&$1$&$0$&$0$&$0$&$0$&$0$&$0$&$0$&$1$\\
$0$&$0$&$0$&$0$&$1$&$1$&$0$&$0$&$0$&$0$&$0$&$1$&$0$\\
$0$&$0$&$0$&$1$&$0$&$0$&$0$&$0$&$0$&$0$&$0$&$1$&$1$\\
$0$&$0$&$0$&$1$&$0$&$1$&$0$&$0$&$0$&$0$&$1$&$0$&$0$\\
$0$&$0$&$0$&$1$&$1$&$0$&$0$&$0$&$0$&$0$&$1$&$0$&$1$\\
$0$&$0$&$0$&$1$&$1$&$1$&$0$&$0$&$0$&$0$&$1$&$1$&$0$\\
$0$&$0$&$1$&$-$&$-$&$0$&$0$&$0$&$0$&$0$&$1$&$1$&$1$\\
$0$&$0$&$1$&$0$&$0$&$1$&$0$&$0$&$0$&$1$&$0$&$0$&$0$\\
$0$&$0$&$1$&$0$&$1$&$1$&$0$&$0$&$0$&$1$&$0$&$1$&$0$\\
$0$&$0$&$1$&$1$&$0$&$1$&$0$&$0$&$0$&$1$&$1$&$0$&$0$\\
$0$&$0$&$1$&$1$&$1$&$1$&$0$&$0$&$0$&$1$&$1$&$1$&$0$\\
$0$&$1$&$0$&$0$&$0$&$0$&$0$&$0$&$0$&$1$&$0$&$0$&$1$\\
$0$&$1$&$0$&$0$&$0$&$1$&$0$&$0$&$1$&$0$&$0$&$0$&$0$\\
$0$&$1$&$0$&$0$&$1$&$0$&$0$&$0$&$1$&$0$&$0$&$0$&$1$\\
$0$&$1$&$0$&$0$&$1$&$1$&$0$&$0$&$1$&$0$&$0$&$1$&$0$\\
$0$&$1$&$0$&$1$&$0$&$0$&$0$&$0$&$1$&$0$&$0$&$1$&$1$\\
$0$&$1$&$0$&$1$&$0$&$1$&$0$&$0$&$1$&$0$&$1$&$0$&$0$\\
$0$&$1$&$0$&$1$&$1$&$0$&$0$&$0$&$1$&$0$&$1$&$0$&$1$\\
$0$&$1$&$0$&$1$&$1$&$1$&$0$&$0$&$1$&$0$&$1$&$1$&$0$\\
$0$&$1$&$1$&$-$&$-$&$0$&$0$&$0$&$1$&$0$&$1$&$1$&$1$\\
$0$&$1$&$1$&$0$&$0$&$1$&$0$&$0$&$1$&$1$&$0$&$0$&$0$\\
$0$&$1$&$1$&$0$&$1$&$1$&$0$&$0$&$1$&$1$&$0$&$1$&$0$\\
$0$&$1$&$1$&$1$&$0$&$1$&$0$&$0$&$1$&$1$&$1$&$0$&$0$\\
$0$&$1$&$1$&$1$&$1$&$1$&$0$&$0$&$1$&$1$&$1$&$1$&$0$\\
$1$&$-$&$0$&$0$&$0$&$0$&$0$&$0$&$1$&$1$&$0$&$0$&$1$\\
$1$&$0$&$0$&$0$&$0$&$1$&$0$&$1$&$0$&$0$&$0$&$0$&$0$\\
$1$&$0$&$0$&$0$&$1$&$0$&$0$&$1$&$0$&$0$&$0$&$0$&$1$\\
$1$&$0$&$0$&$0$&$1$&$1$&$0$&$1$&$0$&$0$&$0$&$1$&$0$\\
$1$&$-$&$0$&$1$&$0$&$0$&$0$&$0$&$1$&$1$&$0$&$1$&$1$\\
$1$&$0$&$0$&$1$&$0$&$1$&$0$&$1$&$0$&$0$&$1$&$0$&$0$\\
$1$&$0$&$0$&$1$&$1$&$0$&$0$&$1$&$0$&$0$&$1$&$0$&$1$\\
$1$&$0$&$0$&$1$&$1$&$1$&$0$&$1$&$0$&$0$&$1$&$1$&$0$\\
$1$&$-$&$1$&$-$&$0$&$0$&$0$&$0$&$1$&$1$&$1$&$1$&$1$\\
$1$&$0$&$1$&$0$&$0$&$1$&$0$&$1$&$0$&$1$&$0$&$0$&$0$\\
$1$&$0$&$1$&$-$&$1$&$0$&$0$&$1$&$0$&$0$&$1$&$1$&$1$\\
$1$&$0$&$1$&$0$&$1$&$1$&$0$&$1$&$0$&$1$&$0$&$1$&$0$\\
$1$&$0$&$1$&$1$&$0$&$1$&$0$&$1$&$0$&$1$&$1$&$0$&$0$\\
$1$&$0$&$1$&$1$&$1$&$1$&$0$&$1$&$0$&$1$&$1$&$1$&$0$\\
$1$&$1$&$0$&$0$&$0$&$1$&$0$&$1$&$1$&$0$&$0$&$0$&$0$\\
$1$&$1$&$0$&$0$&$1$&$0$&$0$&$1$&$1$&$0$&$0$&$0$&$1$\\
$1$&$1$&$0$&$0$&$1$&$1$&$0$&$1$&$1$&$0$&$0$&$1$&$0$\\
$1$&$1$&$0$&$1$&$0$&$1$&$0$&$1$&$1$&$0$&$1$&$0$&$0$\\
$1$&$1$&$0$&$1$&$1$&$0$&$0$&$1$&$1$&$0$&$1$&$0$&$1$\\
$1$&$1$&$0$&$1$&$1$&$1$&$0$&$1$&$1$&$0$&$1$&$1$&$0$\\
$1$&$1$&$1$&$0$&$0$&$1$&$0$&$1$&$1$&$1$&$0$&$0$&$0$\\
$1$&$1$&$1$&$-$&$1$&$0$&$0$&$1$&$1$&$0$&$1$&$1$&$1$\\
$1$&$1$&$1$&$0$&$1$&$1$&$0$&$1$&$1$&$1$&$0$&$1$&$0$\\
$1$&$1$&$1$&$1$&$0$&$1$&$0$&$1$&$1$&$1$&$1$&$0$&$0$\\
$1$&$1$&$1$&$1$&$1$&$1$&$0$&$1$&$1$&$1$&$1$&$1$&$0$\\

\end{tabular}
\end{center}
As expressões booleanas minimizadas são:

$Carry =  \overline{D11}  \cdot  \overline{D10}  \cdot  \overline{D03}  \cdot  \overline{D02}  \cdot  \overline{D01}  \cdot  \overline{D00} $~\\
$o11 =  \overline{D11}  \cdot  \overline{D10}  \cdot  \overline{D03}  \cdot  \overline{D02}  \cdot  \overline{D01}  \cdot  \overline{D00} +D11 \cdot D00+D11 \cdot D01$~\\
$o10 = D10 \cdot D00+D10 \cdot D01+D10 \cdot D02+D10 \cdot D03+D11 \cdot  \overline{D01}  \cdot  \overline{D00} $~\\
$o03 = D03 \cdot D00+D10 \cdot  \overline{D03}  \cdot  \overline{D02}  \cdot  \overline{D01}  \cdot  \overline{D00} +D11 \cdot  \overline{D01}  \cdot  \overline{D00} $~\\
$o02 = D02 \cdot D00+D02 \cdot D01+D03 \cdot  \overline{D00} $~\\
$o01 =  \overline{D11}  \cdot  \overline{D10}  \cdot  \overline{D01}  \cdot  \overline{D00} +D01 \cdot D00+D02 \cdot  \overline{D01}  \cdot  \overline{D00} +D03 \cdot  \overline{D00} $~\\
$o00 =  \overline{D00} $~\\

Neste caso, o sinal $Carry$ indica underflow, ou seja, a transição de $000$ para $101$.

As Figuras~\ref{fig:decmod24.png} e~\ref{fig:decmod24-red.png} mostram a implementação lógica e sua versão em redstone.

\imgbig{decmod24.png}{Circuito lógico para decremento em módulo 24.}
\imgbig{decmod24-red.png}{Implementação em redstone do circuito de decremento em módulo 24.}

\paragraph{Considerações de Projeto}

Embora funcional, o contador apresentado não é autocorretivo, pois estados inválidos não são redirecionados explicitamente para o ciclo válido. Em aplicações físicas sujeitas a ruído ou falhas transitórias, isso pode levar a comportamentos inesperados.

Uma alternativa mais robusta consiste em tratar os estados inválidos como condições de correção, forçando a transição para um estado válido (tipicamente $000$), ou utilizando tais estados como \textit{don't care} durante a minimização lógica para otimizar o circuito.
